\documentclass[12pt]{article}
\usepackage[spanish]{babel}
\usepackage[utf8]{inputenc}
\usepackage[T1]{fontenc}
\usepackage{geometry}
\usepackage{graphicx}
\usepackage{setspace}
\usepackage{titlesec}
\usepackage{listings}
\usepackage{xcolor}
\usepackage{float}
\usepackage[export]{adjustbox}

\geometry{margin=1in}
\onehalfspacing

\definecolor{codegray}{rgb}{0.95,0.95,0.95}

\lstset{
    backgroundcolor=\color{codegray},
    basicstyle=\ttfamily\small,
    breaklines=true,
    frame=single
}

\begin{document}

\begin{titlepage}
    \centering

    {\scshape\Large Programación funcional \par}
    \vspace{1.5cm}

    {\huge\bfseries Base de Conocimientos Familiar en Prolog \par}
    \vspace{0.5cm}
    {\Large Tarea 2 \par}

    \vfill

    {\large
    Nombre del Alumno: Emilio Izquierdo Montero {\hspace{6cm}} \\
    \vspace{0.3cm}
    Profesor: Octavo Guzman Peñaloza {\hspace{6cm}} \\
    \vspace{0.3cm}
    Grupo / Semestre: 8 {\hspace{6cm}} \\
    \vspace{0.3cm}
    Fecha: 29 de mayo del 2026 {\hspace{6cm}}
    \par}

    \vfill

\end{titlepage}

\tableofcontents
\newpage

\section{Introducción}

En esta actividad se desarrolló una base de conocimientos familiar utilizando Prolog. El sistema permite inferir relaciones de parentesco a partir de hechos básicos definidos mediante las relaciones padre/2 y madre/2.

La práctica incluye más de 15 personas, 3 generaciones familiares, diferentes ramas familiares y relaciones tanto paternas como maternas.

\section{Objetivo}

Construir un sistema lógico familiar capaz de inferir relaciones de parentesco utilizando programación lógica en Prolog.

\section{Premisa Fundamental}

La base de conocimientos fue construida únicamente utilizando los hechos básicos:

\begin{lstlisting}[language=Prolog]
padre(Padre, Hijo).
madre(Madre, Hijo).
\end{lstlisting}

A partir de estas relaciones se infirieron automáticamente las demás relaciones familiares mediante reglas.

\section{Estructura de la Base de Conocimientos}

\subsection{Hechos principales}

\begin{lstlisting}[language=Prolog]

/* PADRES */

padre(roberto, carlos).
padre(roberto, patricia).
padre(roberto, miguel).

padre(carlos, daniel).
padre(carlos, sofia).

padre(miguel, andres).
padre(miguel, camila).

padre(jose, laura).
padre(jose, ricardo).

padre(ricardo, valeria).
padre(ricardo, fernando).

padre(daniel, mateo).

padre(andres, lucia).


/* MADRES */

madre(elena, carlos).
madre(elena, patricia).
madre(elena, miguel).

madre(laura, daniel).
madre(laura, sofia).

madre(silvia, andres).
madre(silvia, camila).

madre(marta, laura).
madre(marta, ricardo).

madre(patricia, valeria).
madre(patricia, fernando).

madre(paula, mateo).

madre(karla, lucia).

\end{lstlisting}

\subsection{Captura de los hechos}

\begin{figure}[H]
    \centering
    \includegraphics[width=0.9\textwidth, height=0.4\textheight, keepaspectratio]{/home/emilio/Descargas/hechos_base.png}
    \caption{Captura de pantalla de los hechos principales en Prolog}
    \label{fig:hechos}
    \small\textit{Se muestran las relaciones padre/2 y madre/2 definidas para la familia.}
\end{figure}

\section{Relaciones inferidas}

\subsection{Relaciones básicas}

\begin{itemize}
    \item hijo/2
    \item hija/2
    \item hermano/2
    \item hermana/2
\end{itemize}

\begin{lstlisting}[language=Prolog]

hijo(X,Y) :-
    padre(Y,X).

hijo(X,Y) :-
    madre(Y,X).

hija(X,Y) :-
    padre(Y,X).

hija(X,Y) :-
    madre(Y,X).

hermano(X,Y) :-
    padre(P,X),
    padre(P,Y),
    madre(M,X),
    madre(M,Y),
    X \= Y.

hermana(X,Y) :-
    padre(P,X),
    padre(P,Y),
    madre(M,X),
    madre(M,Y),
    X \= Y.

\end{lstlisting}

\subsection{Captura de relaciones básicas}

\begin{figure}[H]
    \centering
    \includegraphics[width=0.9\textwidth, height=0.4\textheight, keepaspectratio]{/home/emilio/Descargas/relaciones_basicas.png}
    \caption{Captura de las reglas para relaciones básicas (hijo, hija, hermano, hermana)}
    \label{fig:basicas}
    \small\textit{Reglas implementadas en Prolog para inferir parentesco básico.}
\end{figure}

\subsection{Relaciones intermedias}

\begin{itemize}
    \item abuelo/2
    \item abuela/2
    \item tio/2
    \item tia/2
    \item primo/2
    \item prima/2
\end{itemize}

\begin{lstlisting}[language=Prolog]

abuelo(X,Y) :-
    padre(X,Z),
    padre(Z,Y).

abuelo(X,Y) :-
    padre(X,Z),
    madre(Z,Y).

abuela(X,Y) :-
    madre(X,Z),
    padre(Z,Y).

abuela(X,Y) :-
    madre(X,Z),
    madre(Z,Y).

tio(X,Y) :-
    hermano(X,Z),
    padre(Z,Y).

tio(X,Y) :-
    hermano(X,Z),
    madre(Z,Y).

tia(X,Y) :-
    hermana(X,Z),
    padre(Z,Y).

tia(X,Y) :-
    hermana(X,Z),
    madre(Z,Y).

\end{lstlisting}

\subsection{Captura de relaciones intermedias}

\begin{figure}[H]
    \centering
    \includegraphics[width=0.9\textwidth, height=0.4\textheight, keepaspectratio]{/home/emilio/Descargas/relaciones_intermedias.png}
    \caption{Captura de las reglas para relaciones intermedias (abuelo, abuela, tío, tía)}
    \label{fig:intermedias}
    \small\textit{Reglas que combinan relaciones básicas para formar parentescos más complejos.}
\end{figure}

\subsection{Relaciones avanzadas}

\begin{itemize}
    \item ancestro/2
    \item descendiente/2
    \item familiar\_directo/2
\end{itemize}

\begin{lstlisting}[language=Prolog]

ancestro(X,Y) :-
    padre(X,Y).

ancestro(X,Y) :-
    madre(X,Y).

ancestro(X,Y) :-
    padre(X,Z),
    ancestro(Z,Y).

ancestro(X,Y) :-
    madre(X,Z),
    ancestro(Z,Y).

descendiente(X,Y) :-
    ancestro(Y,X).

\end{lstlisting}

\subsection{Captura de relaciones avanzadas}

\begin{figure}[H]
    \centering
    \includegraphics[width=0.9\textwidth, height=0.4\textheight, keepaspectratio]{/home/emilio/Descargas/relaciones_avanzadas.png}
    \caption{Captura de las reglas para relaciones avanzadas (ancestro, descendiente)}
    \label{fig:avanzadas}
    \small\textit{Reglas recursivas que permiten inferir múltiples generaciones.}
\end{figure}

\section{Consultas obligatorias}

\subsection{Consulta de abuelos}

\begin{lstlisting}[language=Prolog]
abuelo(X, mateo).
\end{lstlisting}

Resultado esperado:

\begin{lstlisting}
X = carlos ;
X = roberto.
\end{lstlisting}

\subsection{Captura de la consulta}

\begin{figure}[H]
    \centering
    \includegraphics[width=0.9\textwidth, height=0.4\textheight, keepaspectratio]{/home/emilio/Descargas/consulta1.png}
    \caption{Consulta ejecutada: \texttt{abuelo(X, mateo)}}
    \label{fig:cons_abuelos}
    \small\textit{Resultado: Los abuelos de mateo son carlos y roberto.}
\end{figure}

\subsection{Consulta de hermanos}

\begin{lstlisting}[language=Prolog]
hermano(X, sofia).
\end{lstlisting}

Resultado esperado:

\begin{lstlisting}
X = daniel.
\end{lstlisting}

\subsection{Captura de la consulta}

\begin{figure}[H]
    \centering
    \includegraphics[width=0.9\textwidth, height=0.4\textheight, keepaspectratio]{/home/emilio/Descargas/consultahermanos.png}
    \caption{Consulta ejecutada: \texttt{hermano(X, sofia)}}
    \label{fig:cons_hermanos}
    \small\textit{Resultado: La única hermana de sofia es daniel.}
\end{figure}

\subsection{Consulta de ancestros}

\begin{lstlisting}[language=Prolog]
ancestro(X, lucia).
\end{lstlisting}

Resultado esperado:

\begin{lstlisting}
X = andres ;
X = karla ;
X = miguel ;
X = silvia ;
X = roberto ;
X = elena.
\end{lstlisting}

\subsection{Captura de la consulta}

\begin{figure}[H]
    \centering
    \includegraphics[width=0.9\textwidth, height=0.4\textheight, keepaspectratio]{/home/emilio/Descargas/ancestros.png}
    \caption{Consulta ejecutada: \texttt{ancestro(X, lucia)}}
    \label{fig:cons_ancestros}
    \small\textit{Resultado: Todos los ancestros de lucia en orden de cercanía.}
\end{figure}

\subsection{Consulta de tíos}

\begin{lstlisting}[language=Prolog]
tio(X, mateo).
\end{lstlisting}

Resultado esperado:

\begin{lstlisting}
X = sofia.
\end{lstlisting}

\subsection{Captura de la consulta}

\begin{figure}[H]
    \centering
    \includegraphics[width=0.9\textwidth, height=0.4\textheight, keepaspectratio]{/home/emilio/Descargas/consultatios.png}
    \caption{Consulta ejecutada: \texttt{tio(X, mateo)}}
    \label{fig:cons_tios}
    \small\textit{Resultado: El tío de mateo es sofia.}
\end{figure}

\subsection{Consulta de primos}

\begin{lstlisting}[language=Prolog]
primo(X, lucia).
\end{lstlisting}

Resultado esperado:

\begin{lstlisting}
false.
\end{lstlisting}

\subsection{Captura de la consulta}

\begin{figure}[H]
    \centering
    \includegraphics[width=0.9\textwidth, height=0.4\textheight, keepaspectratio]{/home/emilio/Descargas/lucianotieneprimos.png}
    \caption{Consulta ejecutada: \texttt{primo(X, lucia)}}
    \label{fig:cons_primos}
    \small\textit{Resultado: No hay primos registrados para lucia en la base de conocimientos.}
\end{figure}

\section{Ejemplo conceptual}

\subsection{Hechos}

\begin{lstlisting}[language=Prolog]

padre(juan, carlos).
madre(maria, carlos).

padre(carlos, ana).
madre(laura, ana).

\end{lstlisting}

\subsection{Regla inferida}

\begin{lstlisting}[language=Prolog]

abuelo(X,Y) :-
    padre(X,Z),
    padre(Z,Y).

\end{lstlisting}

\subsection{Consulta}

\begin{lstlisting}[language=Prolog]
abuelo(X, ana).
\end{lstlisting}

\subsection{Resultado esperado}

\begin{lstlisting}
X = juan.
\end{lstlisting}

\section{Conclusión}

Con esta práctica fue posible desarrollar una base de conocimientos familiar utilizando Prolog y programación lógica. A partir de hechos básicos se lograron inferir múltiples relaciones familiares mediante reglas y recursividad.

\end{document}
